\section{La ruta hacia los robots de propósito general: de la Automation a la Superhuman}

Esta sección resume cómo divide Tan Jie (谭杰) las etapas del desarrollo de los robots. Menciona cinco etapas: Automation, Teleoperation, Narrow Generalist, Home Generalist y Superhuman Capability. Esta hoja de ruta ayuda a no mezclar todas las capacidades de los robots en un solo problema.

\begin{figure}[H]
\centering
\includegraphics[width=0.96\textwidth]{robot-generalist-stages.png}
\caption{Etapas de los robots de propósito general: la ruta de largo plazo desde la automatización fija hasta las capacidades sobrehumanas. Diagrama conceptual propio, elaborado a partir de la conversación de 01:17:35--02:06:16.}
\end{figure}

\begin{knowledgebox}{Lectura de la figura: la «generalidad» es distinta en cada etapa}
La Automation consiste en reglas fijas; en la Teleoperation el hardware ya es utilizable, pero una persona lo controla a distancia; el Narrow Generalist generaliza su inteligencia dentro de un dominio acotado; el Home Generalist resuelve múltiples tareas generales en el hogar, y la Superhuman supera a los humanos en ciertas capacidades. No debe confundirse la generalización en un dominio acotado con la generalidad en el hogar.
\end{knowledgebox}

\subsection{Specialist y Generalist}

Las primeras implementaciones prácticas probablemente seguirán ocurriendo en escenarios verticales como la manufactura, la logística, los supermercados, el doblado de ropa o las servilletas. Para un specialist vertical es más fácil encontrar demanda, ROI y entornos controlables. Sin embargo, si llega a consolidarse un verdadero generalist, podrá hacer las tareas de un specialist y muchas más, y muchos specialists verán reducido su espacio. El problema es que un generalist requiere más tiempo y más datos.

\begin{warningbox}{No hay que descartar la comercialización intermedia en nombre de la forma final}
El robot humanoide de propósito general puede ser una meta de largo plazo, pero en el corto plazo los escenarios verticales todavía pueden generar valor. La ruta especializada y la ruta general no son mutuamente excluyentes; es posible que avancen en paralelo durante muchos años.
\end{warningbox}

\subsection{La seguridad no es un asunto menor}

Tan Jie afirma con claridad que hay que atender la AI safety y la robot safety. Cuando la IA o los robots puedan iterar sobre sí mismos, la humanidad enfrentará un problema de supervivencia. Google DeepMind tiene un Responsibility and Safety Council que revisa el impacto de los modelos y los robots en la sociedad, así como sus consecuencias en materia de seguridad. En el peor de los casos, si las capacidades de los robots rebasan la comprensión que se tiene de su seguridad, habría que detener la expansión de capacidades para que la investigación en seguridad se ponga al día.

\begin{importantbox}{Principios de seguridad en robótica}
Las capacidades y la seguridad deben avanzar a la par. En cada etapa del desarrollo debe hacerse la investigación de seguridad correspondiente; si el avance de las capacidades rebasa la comprensión de la seguridad, debe pausarse la expansión de capacidades para que la safety se ponga al día (catch up).
\end{importantbox}

\subsection{Por qué es fácil sobrestimar a los robots}

Esta subsección añade la advertencia de Tan Jie sobre el ánimo de la industria. Es muy fácil sobrestimar a los robots porque lo que el público ve suele ser la mejor demo: el equipo quizá la grabó diez veces y publicó la mejor toma; el video no muestra cuántas veces falló, los supuestos sobre el entorno, los reinicios manuales, los límites de la tarea ni las intervenciones de seguridad. El espectador tiende a confundir la «mejor muestra» con una «capacidad estable» y, por ello, a creer que el año próximo ya podrá comprar un robot humanoide doméstico.

\begin{warningbox}{Los límites de un video de demo como evidencia}
Un video de un robot puede demostrar que «esto ocurrió bajo ciertas condiciones», pero no que el robot tenga una capacidad desplegable. Para juzgar si está listo para usarse en la práctica hay que examinar la tasa de éxito en intentos repetidos, la recuperación ante fallas, los cambios de entorno, la diversidad de tareas, el costo, el mantenimiento y los límites de seguridad.
\end{warningbox}

La actitud de Tan Jie es de entusiasmo, pero con cabeza fría: los robots sí se están acelerando y ya empiezan a aparecer escenarios de aplicación práctica; aun así, los «robots humanoid que realmente pueden trabajar» siguen siendo hoy un desierto. Lo que de verdad hay que evitar es cometer dos errores a la vez: sobrestimar las capacidades de corto plazo y subestimar el impacto de largo plazo.

\begin{importantbox}{La percepción correcta de los plazos}
En el corto plazo, los robots todavía están lejos de la generalidad en el hogar; en el mediano plazo, es probable que primero se implementen en escenarios verticales; en el largo plazo, una vez que se consolide de verdad un generalist, el espacio de supervivencia de los specialists se reducirá.
\end{importantbox}

\subsection{Resumen de la sección}

La ruta de los robots debe verse por etapas. En el corto plazo están la generalización en dominios acotados y la implementación vertical; en el mediano plazo, el generalist doméstico, y solo en el largo plazo llegan las capacidades sobrehumanas. Los problemas de seguridad deben incorporarse al proceso de investigación desde las primeras etapas.
